\documentclass[11pt,a4paper]{article}
\usepackage[top=1.7cm,bottom=1.5cm,left=1.5cm,right=1.5cm,headsep=8pt,footskip=16pt]{geometry}
\usepackage{enumitem}
\usepackage{amssymb}
\usepackage{array}
\usepackage[table]{xcolor}
\usepackage{fancyhdr}
\usepackage{lastpage}
\usepackage{multicol}
\usepackage{polyglossia}
\setdefaultlanguage[calendar=gregorian, numerals=maghrib]{arabic}
\setotherlanguage{english}
\IfFontExistsTF{Amiri}{%
  \newfontfamily\arabicfont[Script=Arabic, Scale=1.08]{Amiri}
}{%
  \newfontfamily\arabicfont[Script=Arabic, Scale=1.08, Path=../fonts/, Extension=.ttf,
    UprightFont=*-Regular, BoldFont=*-Bold, ItalicFont=*-Italic,
    BoldItalicFont=*-BoldItalic]{Amiri}
}
\IfFontExistsTF{AlBattar}{%
  \newfontfamily\khatt[Script=Arabic, Scale=1.15]{AlBattar}
}{%
  \newfontfamily\khatt[Script=Arabic, Scale=1.15, Path=../fonts/,
    BoldFont=ae_AlBattar.ttf]{ae_AlBattar.ttf}
}
\renewcommand{\headrulewidth}{0.8pt}
\renewcommand{\footrulewidth}{0.8pt}
\fancyhf{}
\rhead{\khatt المدرسة العليا للأساتذة ببوزريعة}
\lhead{\khatt ملحقة بومرداس}
\lfoot{\khatt الأستاذ: د. بن عبد الرزاق يوسف}
\cfoot{\thepage/\pageref{LastPage}}
\pagestyle{fancy}
\parindent 0pt
\setlength{\columnsep}{0.6cm}
\setlength{\columnseprule}{0.4pt}
\setlist[enumerate,1]{label=\textbf{\arabic*.}, leftmargin=*, itemsep=5pt, topsep=2pt}
\begin{document}

\linespread{1.08}\selectfont
\begin{center}
{\LARGE\khatt\textbf{الاختبار العاشر في الذكاء الاصطناعي}}\\[4pt]
{\normalsize وضعيات إشكالية ــ أسئلة متعددة الاختيارات ــ الفصول من 1 إلى 6}
\end{center}
\vspace{-2pt}\rule{\textwidth}{0.6pt}\\[2pt]
\textbf{التعليمة:} يتضمن الاختبار خمس وضعيات إشكالية. اقرأ كل وضعية بعناية، ثم اختر التحليل أو التصرف الأنسب
من بين الاختيارات الخمسة؛ لكل وضعية إجابة صحيحة واحدة فقط.
اكتب \textbf{حرف} الإجابة الصحيحة في جدول الإجابات في آخر الاختبار.
لكل وضعية 4 نقاط (المجموع 20 نقطة).
\vspace{4pt}

\noindent
{\large\khatt\textbf{الوضعية 1}}\hfill{\small(الفصل 1)}\par\nopagebreak\vspace{3pt}
\setlength{\fboxsep}{7pt}\colorbox{gray!10}{\parbox{\dimexpr\linewidth-14pt}{سألت طالبة أداة ذكاء اصطناعي عن حكم تقديم الخبر على المبتدأ، فقدّمت شرحًا منظّمًا مع قول منسوب إلى سيبويه بين علامتي تنصيص، وإحالة إلى «الكتاب، الجزء الثاني، صفحة 87». بحثت الطالبة في الموضع المذكور من طبعتين فلم تجد القول. وتريد استعمال القول في مذكرة تخرجها. ما الاستنتاج الأدق؟}}\par\nopagebreak
\begin{itemize}[leftmargin=1.0cm, labelsep=6pt, itemsep=3pt, topsep=4pt, parsep=0pt]
  \item[\textbf{أ)}] قد يكون القول مختلقًا أو محرّفًا أو في موضع آخر، لأن الأداة تولّد إحالات تبدو دقيقة دون استرجاع فعلي؛ فلا تنسبه إلى سيبويه حتى تجده بنصه، وتتحقق من الشرح في مصادر النحو.
  \item[\textbf{ب)}] دقة الإحالة بالجزء والصفحة دليل على أن الأداة رجعت إلى الكتاب فعلًا، والطالبة استعملت طبعات مختلفة، فتعتمد القول والإحالة كما قدّمتهما الأداة دون بحث إضافي، ولا داعي للبحث في طبعات أخرى.
  \item[\textbf{ج)}] القول صحيح، والخطأ في رقم الصفحة فقط، لأن الأداة تحفظ النصوص بدقة وتخطئ في الأرقام وحدها، فتنسب القول إلى سيبويه وتحيل إلى الكتاب دون تحديد الصفحة، ثم تذكر في الهامش أن الصفحة غير محددة.
  \item[\textbf{د)}] الشرح كله خاطئ بالضرورة لأن الإحالة خاطئة، فترفض قاعدة تقديم الخبر على المبتدأ نفسها، وتبحث عن قاعدة مختلفة، لأن خطأ الإحالة يُبطل كل ما ورد في الإجابة.
  \item[\textbf{هـ)}] تطلب من الأداة تصحيح رقم الصفحة، فإن قدّمت رقمًا آخر اعتمدته، لأن الأداة تراجع نفسها عند الطلب، والإحالة المصحَّحة تصبح موثوقة بعد المراجعة الثانية.
\end{itemize}
\par\vspace{6pt}
\noindent
{\large\khatt\textbf{الوضعية 2}}\hfill{\small(الفصل 2)}\par\nopagebreak\vspace{3pt}
\setlength{\fboxsep}{7pt}\colorbox{gray!10}{\parbox{\dimexpr\linewidth-14pt}{أعلن فريق طلابي أن نموذجه يتعرّف على بحر البيت الشعري بدقة 99\%. وتبيّن أنهم قسموا الأبيات عشوائيًا إلى تدريب واختبار، فوقعت أبيات من القصيدة الواحدة في الجزأين. وعند تجربته على قصائد جديدة تمامًا انخفضت دقته كثيرًا. ويطالب أعضاء الفريق بنشر النتيجة في المجلة الطلابية. ما التقييم الأدق لهذا الإعلان؟}}\par\nopagebreak
\begin{itemize}[leftmargin=1.0cm, labelsep=6pt, itemsep=3pt, topsep=4pt, parsep=0pt]
  \item[\textbf{أ)}] انخفاض الدقة على القصائد الجديدة طبيعي في كل النماذج ولا يدل على خلل في التقييم، فيُعتمد رقم 99\% في التقرير، لأنه يمثّل أفضل أداء ممكن للنموذج في الظروف المثالية.
  \item[\textbf{ب)}] الدقة مضخّمة، والحل زيادة حجم جزء الاختبار من الأبيات المختلطة نفسها إلى النصف، لأن كبر جزء الاختبار يلغي أثر تشابه أبيات القصيدة الواحدة ويعطي دقة واقعية.
  \item[\textbf{ج)}] الدقة المعلنة مضخّمة، لأن أبيات القصيدة الواحدة تتشابه وزنًا وقافية وألفاظًا، فاستفاد النموذج من هذا التسرّب؛ والتقييم السليم يفصل القصائد كاملة بين التدريب والاختبار.
  \item[\textbf{د)}] الدقة مضخّمة لأن النموذج حفظ البحور الستة عشر كلها، والحل حذف بعض البحور من التدريب، لأن النموذج الذي لا يعرف كل البحور يتعلم التعميم أفضل على القصائد الجديدة.
  \item[\textbf{هـ)}] الدقة المعلنة صحيحة، لأن التقسيم العشوائي هو أعدل طريقة دائمًا، وانخفاض الدقة على القصائد الجديدة سببه أنها رديئة الوزن، فالنموذج لا يخطئ إلا مع الشعر المكسور.
\end{itemize}
\par\vspace{6pt}
\noindent
{\large\khatt\textbf{الوضعية 3}}\hfill{\small(الفصل 3)}\par\nopagebreak\vspace{3pt}
\setlength{\fboxsep}{7pt}\colorbox{gray!10}{\parbox{\dimexpr\linewidth-14pt}{طلب أستاذ «خطة علاجية لتلاميذي الضعاف في الإملاء»، فجاءت خطة عامة تصلح لأي قسم. لديه نتائج تشخيص تبيّن أن أغلب الأخطاء في الهمزة المتوسطة والتاء المربوطة، وأن الحصة المتاحة ساعة أسبوعيًا لمدة ستة أسابيع. ولا يملك وقتًا لإعداد الخطة بنفسه من البداية. ما الطلب الأنجع للحصول على خطة مناسبة؟}}\par\nopagebreak
\begin{itemize}[leftmargin=1.0cm, labelsep=6pt, itemsep=3pt, topsep=4pt, parsep=0pt]
  \item[\textbf{أ)}] يزوّد الأداة بأسماء التلاميذ وعلاماتهم الكاملة في كل المواد، لأن كثرة البيانات الشخصية تجعل الخطة أدق، ولا حاجة إلى تحديد نوع الأخطاء أو مدة الحصة أو عدد الأسابيع، فتتمكن الأداة من معرفة كل تلميذ معرفة شخصية.
  \item[\textbf{ب)}] يزوّد الأداة بنوع الأخطاء فقط، ويطلب خطة لكل قواعد الإملاء في العربية، لأن الخطة الشاملة أنفع من الخطة المركّزة، والتلاميذ الضعاف يحتاجون إلى مراجعة كل القواعد معًا، فيكون العلاج شاملًا ومتكاملًا لكل التلاميذ.
  \item[\textbf{ج)}] يطلب عشر خطط عامة مختلفة ويختار أطولها، لأن الخطة الأطول تتضمن بالضرورة ما يناسب أخطاء تلاميذه، ولا حاجة إلى تزويد الأداة بأي معلومات عن القسم أو التشخيص.
  \item[\textbf{د)}] يزوّد الأداة بنتائج التشخيص وعدد التلاميذ ومدة الحصة وعدد الأسابيع، ويطلب خطة مرحلية تبدأ بالخطأ الأكثر شيوعًا، مع نشاط وتقويم قصير لكل حصة، ويطلب منها أن تسأله عمّا ينقصها قبل البدء.
  \item[\textbf{هـ)}] يعيد الطلب نفسه مضيفًا «خطة دقيقة ومفصّلة جدًا»، لأن طلب التفصيل يجعل الأداة تستنتج أخطاء التلاميذ ومدة الحصص بنفسها، دون حاجة إلى تزويدها بنتائج التشخيص.
\end{itemize}
\par\vspace{6pt}
\noindent
{\large\khatt\textbf{الوضعية 4}}\hfill{\small(الفصل 4)}\par\nopagebreak\vspace{3pt}
\setlength{\fboxsep}{7pt}\colorbox{gray!10}{\parbox{\dimexpr\linewidth-14pt}{سمح أستاذ باستعمال الذكاء الاصطناعي في بحث منزلي، فتبيّن أن بعض الطلبة يملكون اشتراكات مدفوعة في أدوات متقدمة، وآخرين لا يملكون إلا هاتفًا بسيطًا واتصالًا ضعيفًا. وجاءت بحوث الفئة الأولى أكثر إتقانًا في الشكل والتنظيم. واحتجّ طلبة الفئة الثانية على النتائج، وطالبوا بإعادة النظر في طريقة التقويم. ما المعالجة الأعدل؟}}\par\nopagebreak
\begin{itemize}[leftmargin=1.0cm, labelsep=6pt, itemsep=3pt, topsep=4pt, parsep=0pt]
  \item[\textbf{أ)}] يضيف نقاطًا تعويضية ثابتة لكل طالب لا يملك اشتراكًا مدفوعًا، دون تغيير معايير التقويم، لأن التعويض العددي يحقق العدل تلقائيًا مهما اختلفت جودة البحوث المقدّمة، ويضمن رضا الطلبة جميعًا.
  \item[\textbf{ب)}] يمنع الطلبة الذين يملكون اشتراكات مدفوعة من استعمالها، ويسمح للآخرين بالأدوات المجانية فقط، لأن تقييد الفئة المتقدمة وحدها هو الطريقة الوحيدة لتحقيق العدل في البحوث المنزلية.
  \item[\textbf{ج)}] يُلزم كل الطلبة بشراء الاشتراك المدفوع نفسه، لأن توحيد الأداة يحقق تكافؤ الفرص، ويُبقي التقويم على الإتقان الشكلي للبحث بوصفه أوضح معيار للحكم على جودة العمل.
  \item[\textbf{د)}] يمنح علامات أعلى للبحوث الأكثر إتقانًا، لأن امتلاك أدوات متقدمة جزء من اجتهاد الطالب، وتكافؤ الفرص في الوسائل ليس من مسؤولية الأستاذ في التقويم الجامعي.
  \item[\textbf{هـ)}] يضمن وصولًا متكافئًا، كتوفير قاعة حواسيب وأداة مجانية موحّدة، ويبني التقويم على ما يقيس كفاءة الطالب نفسه، كالتعليل والمناقشة الشفهية وتوثيق المصادر، لا على الإتقان الشكلي.
\end{itemize}
\par\vspace{6pt}
\noindent
{\large\khatt\textbf{الوضعية 5}}\hfill{\small(الفصل 6)}\par\nopagebreak\vspace{3pt}
\setlength{\fboxsep}{7pt}\colorbox{gray!10}{\parbox{\dimexpr\linewidth-14pt}{ولّدت أستاذة صورًا لبطاقات مفردات للسنة الأولى ابتدائي: «الكسكس، القندورة، الطاجين، البرنوس». جاءت بعض الصور لأطباق ولباس من ثقافات أخرى تحمل أسماء مشابهة، فاختلط الأمر على زميلة جرّبت البطاقات. وتريد أن يربط الطفل الكلمة بما يراه في محيطه. وستُستعمل البطاقات في أقسام عدة طوال السنة. ما التصرف الأنسب؟}}\par\nopagebreak
\begin{itemize}[leftmargin=1.0cm, labelsep=6pt, itemsep=3pt, topsep=4pt, parsep=0pt]
  \item[\textbf{أ)}] تكتب الكلمات في الطلب بالإنجليزية، لأن الأداة تفهم الإنجليزية أفضل، فتأتي الصور مطابقة للمرجع الجزائري بالضرورة، دون حاجة إلى وصف بصري أو مراجعة لاحقة للبطاقات.
  \item[\textbf{ب)}] تضيف إلى الطلب سياقًا ثقافيًا ووصفًا بصريًا دقيقًا للشيء كما هو في الجزائر، وتراجع كل بطاقة مع زميلة أو مرجع، وتستبدل الصورة غير المطابقة بصورة حقيقية عند الحاجة.
  \item[\textbf{ج)}] تعتمد الصور كما هي، لأن الطفل في السنة الأولى لا يميّز الفروق الثقافية الدقيقة، والمهم أن يحفظ الكلمة مكتوبة، أما مطابقة الصورة لمرجعها في محيطه فأمر ثانوي في هذه المرحلة.
  \item[\textbf{د)}] تضيف السياق الثقافي إلى الطلب، ثم تعتمد الصور الجديدة دون مراجعة، لأن ذكر الجزائر في الطلب يضمن وحده تطابق كل صورة مع مرجعها الثقافي في كل البطاقات.
  \item[\textbf{هـ)}] تتخلى عن الصور وتكتفي ببطاقات مكتوبة، لأن أدوات توليد الصور لا تعرف شيئًا عن الثقافة الجزائرية، ولا يمكن تحسين نتائجها بأي وصف مهما كان دقيقًا ومفصّلًا.
\end{itemize}
\par\vspace{6pt}
\begin{center}
{\khatt\textbf{جدول الإجابات}}\\[3pt]
\renewcommand{\arraystretch}{1.6}
\begin{tabular}{|>{\centering\arraybackslash}m{1.9cm}|*{5}{>{\centering\arraybackslash}m{1.2cm}|}}
\hline
\textbf{رقم السؤال} & \textbf{1} & \textbf{2} & \textbf{3} & \textbf{4} & \textbf{5} \\ \hline
\textbf{الإجابة} &  &  &  &  &  \\[8pt] \hline
\end{tabular}

\vspace{8pt}
\renewcommand{\arraystretch}{2.2}
\begin{tabular}{|>{\centering\arraybackslash}m{7.5cm}|>{\centering\arraybackslash}m{3.2cm}|>{\centering\arraybackslash}m{3.2cm}|}
\hline
\textbf{الاسم واللقب} & \textbf{الفوج} & \textbf{العلامة} \\ \hline
 & & \hfill /20 \\ \hline
\end{tabular}
\end{center}
\end{document}
